\documentclass[11pt]{article}
\usepackage[T1]{fontenc}
\usepackage{lmodern,microtype,amsmath,amssymb,amsthm,mathtools,mathrsfs,tikz-cd}
\usepackage[margin=1in]{geometry}
\usepackage{enumitem,booktabs,longtable,xurl,hyperref}
\hypersetup{colorlinks=true,linkcolor=blue,citecolor=blue,urlcolor=blue,
 pdftitle={Domains, Hoare profunctors, and essential incompleteness},pdfauthor={Florian Lengyel}}
\newtheorem{theorem}{Theorem}[section]
\newtheorem{proposition}[theorem]{Proposition}
\newtheorem{lemma}[theorem]{Lemma}
\newtheorem{corollary}[theorem]{Corollary}
\theoremstyle{definition}
\newtheorem{definition}[theorem]{Definition}
\newtheorem{example}[theorem]{Example}
\theoremstyle{remark}
\newtheorem{remark}[theorem]{Remark}
\newtheorem{problem}[theorem]{Problem}
\newcommand{\R}{\mathscr R}
\newcommand{\D}{\mathsf D}
\newcommand{\Dom}{\operatorname{Dom}}
\newcommand{\ran}{\operatorname{ran}}
\newcommand{\Tot}{\operatorname{Tot}}
\newcommand{\Pos}{\mathbf{Pos}}
\newcommand{\Set}{\mathbf{Set}}
\newcommand{\Prof}{\mathbf{Prof}}
\newcommand{\CE}{\operatorname{CE}}
\newcommand{\N}{\mathbb N}
\newcommand{\op}{\mathrm{op}}
\newcommand{\id}{\operatorname{id}}
\newcommand{\Cons}{\mathsf{Cons}}
\newcommand{\Th}{\mathsf{Th}}
\newcommand{\BDLat}{\mathbf{BDLat}}
\title{Domains, Hoare profunctors,\\and essential incompleteness}
\author{Florian Lengyel}
\date{29 October 2026}
\begin{document}
\maketitle
\begin{abstract}
Restriction domains and ranges determine profunctors for total and partial
correctness. Their composition follows program composition, while predicate
inverse image is contravariant. Partial correctness remains expressible when
weakest liberal preconditions are not representable by computational domains:
for a partial recursive map, direct image on computably enumerable domains has a right adjoint
exactly when its domain is recursive. Two categorical separation arguments give
incompleteness under explicit hypotheses on representation or finite proofs.
A numerical realization constructs a consistent computably enumerable
classical deductive system every consistent computably enumerable extension
of which is incomplete. The construction relates finite computation
codes to positive and negative provability without assuming soundness
of the extensions. An intrinsic construction of the required deductive data
from an abstract recursion category remains a separate existence problem.
\end{abstract}
\tableofcontents

\section{Computation, assertions, and deduction}
Essential incompleteness concerns consistent effective extensions of a
deductive system. A categorical account must therefore connect computational
structure with assertions and deductions, and explain how consistency and
completeness impose incompatible conditions on computational domains.
Di Paola and Heller formulate this problem through categorical analogues of
recursion-theoretic inseparability and creativity, together with an intrinsic
connection to algebraic formal systems \cite[\S1]{DPH}.

There are three constructions to distinguish. A restriction domain describes
where a partial computation is defined. A Hoare relation describes which
preconditions and postconditions a computation satisfies. A deductive
presentation describes which assertions have finite proofs. The first two
are determined by partial-map structure and ranges. The third is explicit
in the numerical construction, where finite computation traces produce
proofs in a classical deductive category. A separator theorem applies when
these proofs give the requisite positive and negative conclusions.

Predicate-transformer semantics associates programming constructs with
transformations of assertions \cite{Dijkstra,DijkstraScholten}.
The categorical semantics of Hoare relations has an established antecedent
in the work of Arthan, Martin, Mathiesen, and Oliva \cite{ArthanEtAl}.
The constructions below specialize the relation between preconditions and
postconditions to restriction domains. They use the local direct-image and
inverse-image adjunction of range restriction categories
\cite[Proposition~8.2]{CockettGarner} and Heller's stable unions
\cite[\S6]{Heller}. The numerical deductive construction makes the
inseparability argument explicit, including its distinction from semantic
incompleteness of program verification \cite[\S6]{Cook}.

Composition is written $gf=g\circ f$, with $f$ performed first. Posets are
regarded as categories with one arrow $p\to q$ precisely when $p\le q$.
All categories used to index coends are small, or are taken in a fixed
universe in which the indicated coends exist. Effectivity in the numerical
sections refers to ordinary computability on natural-number codes.


\section{Restriction domains and direct images}\label{hp:domains}

A \emph{restriction
category} is a category $\R$ equipped with an operation assigning to each
$f:X\to Y$ an endomorphism $\bar f:X\to X$ such that
\begin{equation}\label{hp:restriction-axioms}
 f\bar f=f,\qquad
 \bar f\bar g=\bar g\bar f,\qquad
 \overline{g\bar f}=\bar g\bar f,\qquad
 \bar h f=f\overline{hf}.
\end{equation}
In the second and third identities $f$ and $g$ have the same source; in
the fourth, $f:X\to Y$ and $h:Y\to Z$. An arrow $f:X\to Y$ is
\emph{total} if $\bar f=1_X$. Total arrows form a subcategory
$\Tot(\R)$ with the same objects as $\R$. We assume local smallness and
zero arrows satisfying $\overline{0_{X,Y}}=0_{X,X}$. The restriction
structure is the part of a recursion category used in this section
\cite{CockettLack}.

A \emph{restriction idempotent} on $X$ is an endomorphism $p:X\to X$
with $\bar p=p$. It expresses a domain on $X$. Let $\Dom(X)$ be the set
of these idempotents, ordered by
\[
 p\le q\quad\Longleftrightarrow\quad qp=p.
\]
The meet is $p\wedge q=pq$, the greatest element is $1_X$, and the
least element is $0_X$. Regard this poset as a category: its hom-set
$\Dom(X)(p,q)$ is a singleton if $p\le q$ and is empty otherwise.

Define the contravariant functor
\begin{equation}\label{hp:domain-functor}
 \D:\R^{\op}\longrightarrow\Pos.
\end{equation}
On objects put $\D(X)=\Dom(X)$. For $f:X\to Y$, the monotone map
$\D(f^{\op}):\D(Y)\to\D(X)$ is denoted $f^*$ and is given by
\[
 f^*(q)=\overline{qf}.
\]
The restriction identities imply
\begin{equation}\label{hp:pullback-laws}
 (gf)^*=f^*g^*,\qquad 1_X^*=1_{\D(X)},\qquad
 f^*(pq)=f^*(p)f^*(q),\qquad
 f^*(1_Y)=\bar f,\qquad f^*(0_Y)=0_X.
\end{equation}
These are the identity, composition, and meet laws of the fundamental
domain construction; see \cite[Proposition~3.1]{CockettGarner}.
In particular, inverse image preserves the greatest element precisely
for total arrows.

Suppose that $\R$ has a \emph{range operator}. Its value on $f:X\to Y$
is a restriction idempotent $\ran(f):Y\to Y$ satisfying
\begin{equation}\label{hp:range-axioms}
 \ran(f)f=f,\qquad
 \ran(ef)=e\ran(f),\qquad
 \ran(g\ran(f))=\ran(gf),
\end{equation}
where $e$ is a restriction idempotent on $Y$ and $g:Y\to Z$.
These are the range-restriction axioms, with the assertion that
$\ran(f)$ is a restriction idempotent included in the definition
\cite[Section~8]{CockettGarner}. They give the following universal
property among domains:
\begin{equation}\label{hp:range-least}
 \ran(f)\le e\quad\Longleftrightarrow\quad ef=f.
\end{equation}
Indeed, the forward implication follows from $\ran(f)f=f$, and the
reverse implication follows from $\ran(f)=\ran(ef)=e\ran(f)$.

Define the covariant functor
\[
 \D_!:\R\longrightarrow\Pos
\]
by $\D_!(X)=\D(X)$ and
\[
 \D_!(f)=f_!: \D(X)\longrightarrow\D(Y),\qquad
 f_!(p)=\ran(fp).
\]
If $p\le p'$, every domain containing the range of $fp'$ also contains
the range of $fp$, so \eqref{hp:range-least} proves monotonicity.
Since $\ran(p)=p$ for a restriction idempotent, and by the last range
axiom,
\begin{equation}\label{hp:image-functor}
 (1_X)_!=1_{\D(X)},\qquad (gf)_!=g_!f_!.
\end{equation}

\begin{proposition}[Local adjunction]\label{hp:local-adjunction}
Let $f:X\to Y$ be an arrow in a range restriction category. Restricting
the domain of $f_!$ and the codomain of $f^*$ gives an adjunction
\[
 f_!:\D(X)\mathbin{\downarrow}\bar f
 \ \rightleftarrows\ \D(Y):f^*,\qquad f_!\dashv f^*,
\]
where $\D(X)\mathbin{\downarrow}\bar f=\{p:p\le\bar f\}$.
For arbitrary $p\in\D(X)$ and $q\in\D(Y)$,
\begin{equation}\label{hp:partial-adjunction}
 f_!(p)\le q\quad\Longleftrightarrow\quad
 p\bar f\le f^*(q).
\end{equation}
Moreover, direct image satisfies Frobenius reciprocity:
\begin{equation}\label{hp:frobenius}
 f_!\bigl(p\wedge f^*(q)\bigr)=f_!(p)\wedge q.
\end{equation}
\end{proposition}
\begin{proof}
First suppose $p\le\bar f$. By \eqref{hp:range-least},
$f_!(p)\le q$ is equivalent to $qfp=fp$. If $p\le\overline{qf}$,
the last restriction identity gives $qfp=f\overline{qf}p=fp$.
Conversely, taking restrictions of $qfp=fp$ gives
$\overline{qf}p=p$, since $p\le\bar f$. This proves the adjunction.
For arbitrary $p$, use $fp=f\bar f p$ to obtain
\eqref{hp:partial-adjunction}. Finally,
$fp\overline{qf}=qfp$, so
\[
 f_!\bigl(p\wedge f^*(q)\bigr)
 =\ran(qfp)=q\ran(fp).
\]
This is also the local adjunction of
\cite[Proposition~8.2]{CockettGarner}.
\end{proof}

\section{Hoare assertions as profunctors}\label{hp:profunctors}

For small categories $A,B$, a \emph{profunctor} $P:A\nrightarrow B$
means a functor
\[
 P:A^{\op}\times B\longrightarrow\Set.
\]
For $P:A\nrightarrow B$ and $Q:B\nrightarrow C$, their composite is
\[
 (Q\odot P)(a,c)=\int^{b\in B}P(a,b)\times Q(b,c).
\]
Here the coend is the quotient of the disjoint union of the displayed
products by the identifications induced by arrows of $B$: for
$u:b\to b'$, $x\in P(a,b)$, and $y\in Q(b',c)$,
\[
 [P(1,u)x,y]=[x,Q(u^{\op},1)y].
\]
The identity profunctor on $A$ is its hom-functor $A(-,-)$.
Categories, profunctors, and natural transformations form the bicategory
$\Prof$; its associativity constraints are the canonical reassociations
of these coends. Each $\D(X)$ above is small because $\R$ is locally
small.

For an arrow $f:X\to Y$, define the functors
\begin{equation}\label{hp:hoare-types}
 H_f,K_f:\D(X)^{\op}\times\D(Y)\longrightarrow\Set
\end{equation}
by
\begin{equation}\label{hp:hoare-definitions}
 H_f(p,q)=\D(X)(p,f^*q),\qquad
 K_f(p,q)=\D(Y)(f_!p,q).
\end{equation}
Their values are empty sets or singletons. On morphisms, both functors
use strengthening of preconditions and weakening of postconditions:
if $p'\le p$ and $q\le q'$, the unique maps
\[
 H_f(p,q)\longrightarrow H_f(p',q'),\qquad
 K_f(p,q)\longrightarrow K_f(p',q')
\]
exist by monotonicity. This specifies their functorial action.

In the category of sets and partial functions, $H_f(p,q)$ is inhabited
exactly when $f$ terminates on every input in $p$ and every resulting
output lies in $q$. The value $K_f(p,q)$ is inhabited exactly when every
terminating computation from $p$ has its output in $q$. These are,
respectively, total and partial correctness. The definitions in
\eqref{hp:hoare-definitions} express the same distinction using only
domains, inverse image, and range.

\begin{proposition}\label{hp:hoare-composition}
The assignments $X\mapsto\D(X)$, $f\mapsto H_f$ and
$X\mapsto\D(X)$, $f\mapsto K_f$ define covariant pseudofunctors
\[
 H,K:\R\longrightarrow\Prof,
\]
where $\R$ is regarded as a bicategory with only identity $2$-cells.
Their composition constraints are
\[
 H_g\odot H_f\cong H_{gf},\qquad
 K_g\odot K_f\cong K_{gf}.
\]
\end{proposition}
\begin{proof}
For $f:X\to Y$, $g:Y\to Z$, $p\in\D(X)$, and $r\in\D(Z)$,
the co-Yoneda identity gives
\begin{align*}
 (H_g\odot H_f)(p,r)
 &=\int^{q\in\D(Y)}\D(X)(p,f^*q)\times\D(Y)(q,g^*r)\\
 &\cong\D(X)(p,f^*g^*r)=H_{gf}(p,r),\\
 (K_g\odot K_f)(p,r)
 &=\int^{q\in\D(Y)}\D(Y)(f_!p,q)\times\D(Z)(g_!q,r)\\
 &\cong\D(Z)(g_!f_!p,r)=K_{gf}(p,r).
\end{align*}
For completeness, the first isomorphism sends a pair of inequalities
$p\le f^*q$, $q\le g^*r$ to their composite
$p\le f^*g^*r$. Its inverse uses $q=g^*r$; every other factorization
is identified with this one by the coend relation. The second
isomorphism composes $f_!p\le q$ with $g_!q\le r$ after applying
$g_!$ to the first inequality; its inverse uses $q=f_!p$.
Identity arrows give the hom-profunctors. The associativity and unit
identities hold because both routes compose the same inequalities.
\end{proof}

The covariance of these pseudofunctors is compatible with the
contravariance of $\D$: the right-hand argument $f^*q$ of the hom-set in
$H_f$ converts inverse image into a profunctor from $\D(X)$ to
$\D(Y)$.

\begin{proposition}\label{hp:hoare-comparison}
For $f:X\to Y$ and domains $p,q$,
\begin{align}
 K_f(p,q)\ne\varnothing
 &\quad\Longleftrightarrow\quad p\bar f\le f^*q,
 \label{hp:k-equivalence}\\
 H_f(p,q)\ne\varnothing
 &\quad\Longleftrightarrow\quad
 p\le\bar f\ \text{ and }\ K_f(p,q)\ne\varnothing.
 \label{hp:h-equivalence}
\end{align}
There is a canonical natural transformation $H_f\Rightarrow K_f$.
It is an isomorphism if and only if $f$ is total.
\end{proposition}
\begin{proof}
Equation~\eqref{hp:k-equivalence} is
\eqref{hp:partial-adjunction}. Since $f^*q\le\bar f$,
the condition $p\le f^*q$ is equivalent to the two conditions in
\eqref{hp:h-equivalence}. This gives the natural transformation.
If $f$ is total, the two conditions coincide for all $p,q$.
Conversely, $K_f(1_X,1_Y)$ is always inhabited, whereas
$H_f(1_X,1_Y)$ is inhabited exactly when $1_X\le\bar f$.
\end{proof}

For the zero arrow $0_{X,Y}$, $H_0(p,q)$ is inhabited exactly when
$p=0_X$, whereas $K_0(p,q)$ is always inhabited. Thus the distinction
already appears for a nowhere-defined program.

The restriction order on parallel arrows is
$f\le g$ if $f=g\bar f$. It induces opposite orders on the two
profunctor assignments:
\begin{equation}\label{hp:program-order}
 f\le g\quad\Longrightarrow\quad
 H_f\Rightarrow H_g\quad\text{and}\quad K_g\Rightarrow K_f.
\end{equation}
Indeed, $f^*q=\bar f g^*q\le g^*q$, and
$f_!p=\ran(g\bar fp)\le\ran(gp)=g_!p$. Each displayed natural
transformation is the unique map between the corresponding subterminal
values. Thus extension of a partial program enlarges its total-correctness
relation and restricts its partial-correctness relation.

\begin{proposition}[Detection of domains]\label{hp:detect-domains}
For restriction idempotents $e,d$ on $X$,
\[
 H_e(p,1_X)\ne\varnothing\quad\Longleftrightarrow\quad p\le e,
 \qquad
 H_e\Rightarrow H_d\quad\Longleftrightarrow\quad e\le d.
\]
Consequently $H_e\cong H_d$ implies $e=d$, and
\[
 ed=0_X\quad\Longleftrightarrow\quad H_d\odot H_e\cong H_{0_X}.
\]
\end{proposition}
\begin{proof}
The first equivalence follows from $e^*(1_X)=e$.
If $H_e\Rightarrow H_d$, evaluate it at $(e,1_X)$ to obtain $e\le d$.
The converse follows from \eqref{hp:program-order}. The final statement
uses $H_d\odot H_e\cong H_{de}$ and the fact that restriction
idempotents commute.
\end{proof}

For partial recursive maps on $\N$, singleton domains distinguish arrows: $H_f(\{x\},\{y\})$ is inhabited exactly when
$f(x)=y$. Hence $H_f\cong H_g$ implies $f=g$ in that category.
More precisely, $H_f\Rightarrow H_g$ holds exactly when $f\le g$.

\section{Weakest liberal preconditions and representability}
\label{hp:representability}

For each $q\in\D(Y)$, partial correctness determines the presheaf
\[
 K_f(-,q):\D(X)^{\op}\longrightarrow\Set.
\]
A \emph{weakest liberal precondition} for $q$ is a representing domain
$w\in\D(X)$, that is, a natural isomorphism
\begin{equation}\label{hp:wlp-representing}
 K_f(p,q)\cong\D(X)(p,w)
 \qquad(p\in\D(X)).
\end{equation}
Equivalently, $w$ is the greatest domain satisfying $f_!w\le q$.
If such a domain exists for every $q$, it defines a monotone map
$\operatorname{wlp}_f:\D(Y)\to\D(X)$ and an adjunction
\[
 f_!:\D(X)\rightleftarrows\D(Y):\operatorname{wlp}_f.
\]
The domain of this left adjoint is all of $\D(X)$; compare the local
adjunction in Proposition~\ref{hp:local-adjunction}.

Suppose that $\D(X)$ has finite joins distributing over meets, and that
$\bar f$ has a complement $c$, meaning $c\bar f=0_X$ and
$c\vee\bar f=1_X$. Then every weakest liberal precondition exists and
\begin{equation}\label{hp:complemented-wlp}
 \operatorname{wlp}_f(q)=c\vee f^*q.
\end{equation}
Indeed, distributivity gives
$p\bar f\le f^*q$ if and only if $p\le c\vee f^*q$; apply
\eqref{hp:k-equivalence}. Stable finite unions of domains supply the
distributivity used here.

Existence of a weakest liberal precondition for $0_Y$ alone means that
$\bar f$ has a \emph{pseudocomplement}: a greatest domain disjoint from
$\bar f$. It does not imply the complement equation
$\bar f\vee\operatorname{wlp}_f(0_Y)=1_X$ in an arbitrary domain
lattice. For example, in the chain $0<a<1$, the pseudocomplement of
$a$ is $0$, and $a\vee0\ne1$. This is a domain lattice: take the
one-object category whose arrows are $0,a,1$, whose composition is meet,
and whose restriction and range operators both fix every arrow.

\subsection{The numerical obstruction}

A subset of $\N$ is \emph{recursive} if membership is decidable, and
\emph{computably enumerable} (c.e.) if it is the domain of a partial
recursive function. Let $\CE(\N)$ be the poset of c.e. subsets of
$\N$, ordered by inclusion. These are precisely the restriction
domains on $\N$ in the category of partial recursive maps: a c.e.
set $P$ determines the partial identity with domain $P$, and every
restriction idempotent is such a partial identity. For a partial
recursive $f:\N\rightharpoonup\N$, both $f(P)$ and
$\{n:f(n)\text{ is defined and belongs to }Q\}$ are c.e. when
$P,Q$ are c.e. They give the maps $f_!$ and $f^*$ on this poset.

\begin{theorem}\label{hp:numerical-wlp}
For a partial recursive function $f:\N\rightharpoonup\N$, the map
\[
 f_!:\CE(\N)\longrightarrow\CE(\N),\qquad P\longmapsto f(P),
\]
has a right adjoint if and only if $\operatorname{dom}f$ is recursive.
In fact, representability of $K_f(-,\varnothing)$ already implies that
$\operatorname{dom}f$ is recursive.
\end{theorem}
\begin{proof}
Suppose that a c.e. set $W$ represents $K_f(-,\varnothing)$. Then
\[
 P\subseteq W\quad\Longleftrightarrow\quad
 P\cap\operatorname{dom}f=\varnothing
 \qquad(P\in\CE(\N)).
\]
Taking $P=W$ shows that $W$ is disjoint from $\operatorname{dom}f$.
If $n\notin\operatorname{dom}f$, taking the recursive singleton
$P=\{n\}$ gives $n\in W$. Thus
$W=\N\setminus\operatorname{dom}f$. The domain of a partial recursive
function is c.e.; since its complement is also c.e., dovetailing the two
enumerations decides membership in $\operatorname{dom}f$.

Conversely, if $\operatorname{dom}f$ is recursive, then
\[
 W_Q=(\N\setminus\operatorname{dom}f)\cup f^*Q
\]
is c.e. for every c.e. $Q$. For every c.e. $P$,
$f(P)\subseteq Q$ holds exactly when $P\subseteq W_Q$.
Therefore $Q\mapsto W_Q$ is the required right adjoint.
\end{proof}

Thus a partial-correctness presheaf can exist without a representing
computational domain. At the false postcondition its values describe
all c.e. domains of nontermination, even when their union is not c.e.:
\[
 K_f(P,\varnothing)\ne\varnothing
 \quad\Longleftrightarrow\quad P\cap\operatorname{dom}f=\varnothing.
\]
The use of singleton domains in the proof explains why the numerical
necessity statement is stronger than the abstract pseudocomplement
condition.

To distinguish semantic predicates from their effective presentations,
consider also the maps
\[
 \operatorname{wp}_f,\operatorname{wlp}_f:
 \mathcal P(\N)\longrightarrow\mathcal P(\N)
\]
on the full powerset, defined by
\begin{align*}
 \operatorname{wp}_f(Q)
 &=\{n:f(n)\text{ is defined and belongs to }Q\},\\
 \operatorname{wlp}_f(Q)
 &=(\N\setminus\operatorname{dom}f)\cup\operatorname{wp}_f(Q).
\end{align*}
These satisfy
\begin{equation}\label{hp:halting-predicates}
 \operatorname{wp}_f(\N)=\operatorname{dom}f,
 \qquad
 \operatorname{wlp}_f(\varnothing)=\N\setminus\operatorname{dom}f.
\end{equation}
For recursive $Q$, membership in $\operatorname{wp}_f(Q)$ is c.e.,
whereas membership in $\operatorname{wlp}_f(Q)$ is co-c.e.:
its complement is $\operatorname{wp}_f(\N\setminus Q)$. Neither
predicate is generally recursive. If $Q$ is given by an arithmetic
formula, a formula for either transformed predicate can nevertheless be
constructed effectively from a program index, by coding finite
computations. Construction of such a formula does not decide its truth.

There is also no uniform procedure that, from an arbitrary index for a
partial recursive function with recursive domain, returns a c.e. index
for the complement of that domain. For a machine index $e$, let
$f_e(n)$ simulate machine $e$ on input $0$ and return $0$ if that
simulation halts. The domain of $f_e$ is either $\N$ or
$\varnothing$, so it is recursive in every case. A uniform index for
its complement would enumerate $0$ exactly when that simulation does
not halt. This would make the complement of the halting set c.e.
Pointwise existence in Theorem~\ref{hp:numerical-wlp} therefore does not
give uniform synthesis from program indices.

\section{Stable unions of domains}\label{ct:union-section}

Complementary domains and their inverse images enter the separation
arguments. This section assumes binary restriction products and the universal
range property of Di Paola--Heller \cite[Definition~5.1]{DPH}.
Explicitly, for
$f:X\to Y$, its range receives $f$ and satisfies
\begin{equation}\label{ct:universal-range}
 \beta f=\beta'f\quad\Longrightarrow\quad
 \beta\ran(f)=\beta'\ran(f)
 \qquad(\beta,\beta':Y\to Z).
\end{equation}
This is an additional hypothesis on the range structure.

Write $X\boxtimes Y$ for a binary restriction product. Its projections are
total, and the pairing of $f:Z\to X$ and $g:Z\to Y$ is characterized by
\[
 \pi_1\langle f,g\rangle=f\bar g,\qquad
 \pi_2\langle f,g\rangle=g\bar f,\qquad
 \overline{\langle f,g\rangle}=\bar f\bar g.
\]
The induced product of arrows is a bifunctor, with the usual canonical
associativity and symmetry isomorphisms. Finite powers below have fixed
parenthesizations. Total arrows form a category $\Tot(\R)$ with binary
products and diagonals $\Delta_X=\langle1_X,1_X\rangle$.

\begin{definition}[Stable union]\label{ct:stable-union}
Let $(e_i)_{i\in I}$ be a nonempty finite or countable family of domains on
$X$. A domain $e$ is its \emph{union} if $e_i\le e$ for every $i$ and
\[
 (\forall i,\ \beta e_i=\beta'e_i)
 \quad\Longrightarrow\quad \beta e=\beta'e
 \qquad(\beta,\beta':X\to Y).
\]
It is \emph{stable} if $1_W\boxtimes e$ is a union of
$(1_W\boxtimes e_i)_i$ for every object $W$.
\end{definition}

This is Heller's definition \cite[\S6, p.~1262]{Heller}. Taking $\beta$ to
be an upper-bound domain and $\beta'=1_X$ shows that a union is a least
upper bound, hence unique. Write it as $\bigvee_i e_i$ and put
$\bigvee\varnothing=0_X$.

\begin{lemma}\label{ct:union-calculus}
Under \eqref{ct:universal-range}, multiplication by a fixed domain and
direct image preserve unions. Inverse image along a total arrow preserves
stable unions. If $e=\bigvee_k e_k$ is a stable union of an increasing
sequence, then, for every $n\ge1$ and object $P$,
\[
 1_P\boxtimes e^{\boxtimes n}
 =\bigvee_k(1_P\boxtimes e_k^{\boxtimes n}).
\]
\end{lemma}
\begin{proof}
Agreement after all $de_i$ implies agreement after $de$ by applying the
union property to $\beta d,\beta'd$. For direct image, agreement after
all $\ran(fe_i)$ first implies agreement after all $fe_i$, then after
$fe$, and finally after $\ran(fe)$ by \eqref{ct:universal-range}.

For a total $h:X\to Y$, put $g=\langle1_X,h\rangle$ and $r=\ran(g)$.
For every domain $v$ on $Y$, the restriction and range identities give
\[
 (1_X\boxtimes v)r=\ran(g h^*v),\qquad
 (\pi_1)_!((1_X\boxtimes v)r)=h^*v.
\]
Indeed, $(1_X\boxtimes v)g=g h^*v$, and $\pi_1g=1_X$.
If $v=\bigvee_i v_i$ is stable, multiply the union of
$1_X\boxtimes v_i$ by $r$ and apply $(\pi_1)_!$.
The first two assertions give $h^*v=\bigvee_i h^*v_i$.
Applying the same argument to $1_W\boxtimes h$ proves stability of this
union.

Finally, each mixed product
$e_{k_1}\boxtimes\cdots\boxtimes e_{k_n}$ is below
$e_k^{\boxtimes n}$ for $k=\max_j k_j$. Agreement after the latter
therefore implies agreement after every mixed product. Apply stability
successively in the $n$ coordinates, retaining the fixed factor $1_P$.
\end{proof}

Assume finite stable unions exist. Each $\D(X)$ is then a bounded
distributive lattice. Let $\BDLat$ denote the category of such lattices
and maps preserving finite meets, finite joins, and both bounds. There is
a functor
\[
 E:\Tot(\R)^{\op}\longrightarrow\BDLat.
\]
It assigns $E(X)=\D(X)$ and $E(f^{\op})=f^*$. The restriction identities
give its functor laws and meet preservation; totality gives $f^*1=1$,
and Lemma~\ref{ct:union-calculus} gives join preservation.


\section{Categorical separation and proof codes}\label{ct:transfer-section}

Retain finite stable unions. A pair of domains $(e_+,e_-)$ on $X$ is
\emph{complementary} if $e_+e_-=0_X$ and $e_+\vee e_-=1_X$.
It \emph{separates} domains $a,b$ if $a\le e_+$ and $b\le e_-$.
Disjoint domains $a,b$ are \emph{inseparable} if no complementary pair
separates them. Below a domain $s$, the relative definitions replace
$1_X$ by $s$. This agrees with the complemented-domain separation
criterion of Di Paola--Heller on the ambient object
\cite[Definition 8.14]{DPH}. Relative inseparability implies ambient
inseparability: intersect an ambient complementary pair with $s$.

On an assertion domain $s$ on $U$, supply an arrow
\[
 \nu:U\to U,\qquad \bar\nu=s,\qquad s\nu=\nu.
\]
For $t\le s$, define $t^-=\nu^*t$. Call $t$ \emph{consistent} if
$tt^-=0_U$, and \emph{complete} if $t\vee t^-=s$.
In syntax, $\nu$ constructs the code of a negated assertion. Its
inverse-image action on domains is monotone.

\begin{theorem}[Incompleteness from inseparable domains]\label{ct:certificate-transfer}
Let $a,b\in\D(X)$ be inseparable. Suppose there are an assertion domain
$s\in\D(U)$, an arrow $\nu$ as above, a proof-conclusion map
$d_0:V\to U$ with $t_0=\ran(d_0)\le s$, a total arrow
$\theta:X\to U$ with $s\theta=\theta$, and arrows $u,v:X\to V$ satisfying
\begin{equation}\label{ct:certificate-equations}
 \bar u=a,\quad d_0u=\theta a,
 \qquad
 \bar v=b,\quad d_0v=\nu\theta b.
\end{equation}
Every consistent domain $t$ with $t_0\le t\le s$ is incomplete.
\end{theorem}
\begin{proof}
Since $t_0d_0=d_0$, the code equations imply
\[
 a\le\theta^*t_0,\qquad
 b\le\theta^*\nu^*t_0.
\]
For example, $t_0\theta a=\theta a$; restricting both sides and using
totality of $\theta$ gives the first inequality. For the second,
$\nu\theta$ is total because $\bar\nu=s$ and $s\theta=\theta$.
Monotonicity then gives
$a\le\theta^*t$ and $b\le\theta^*t^-$.
If $t$ were consistent and complete, total inverse image would make
$(\theta^*t,\theta^*t^-)$ complementary on $X$, since
$\theta^*s=1_X$. This would separate $a,b$.
\end{proof}

In a deductive presentation, $d_0$ may be a map from proof codes
to their conclusions. The equations \eqref{ct:certificate-equations}
then require actual constructions of positive and negative proofs.
Deductive closure specifies the theory domains to which the theorem
applies; the contradiction uses the code equations, consistency,
completeness, and inseparability.

\subsection{A representation condition implying inseparability}

There is a second route in which inseparability follows from a
representation property. Fix a nontrivial object $U$, meaning
$0_U\ne1_U$, a domain $s$ on $U$, disjoint domains $p_+,p_-\le s$,
and a total arrow
\[
 \sigma:U\boxtimes U\to U,\qquad s\sigma=\sigma.
\]
For every total $g:U\to U$, put
$f_g=\sigma(g\boxtimes1_U):U\boxtimes U\to U$.
Assume that for every complementary pair $(e_+,e_-)$ on $U$ there is
such a $g$ with
\begin{equation}\label{ct:representation}
 f_g^*p_+=\pi_2^*e_+,\qquad
 f_g^*p_-=\pi_2^*e_-.
\end{equation}

\begin{theorem}\label{ct:representation-inseparable}
Under these hypotheses, $p_+,p_-$ are inseparable relative to $s$.
\end{theorem}
\begin{proof}
Suppose a complementary pair $(d_+,d_-)$ relative to $s$ separates them.
Put $\delta=\sigma\Delta_U$. Total inverse image makes
$e_+=\delta^*d_-$ and $e_-=\delta^*d_+$ complementary on $U$.
Choose $g$ satisfying \eqref{ct:representation}, and put
$\gamma=\langle1_U,g\rangle$ and $h=\delta g$. Then
\[
 f_g\gamma=h,\qquad\pi_2\gamma=g.
\]
The first identity gives the commutative square
\[
\begin{tikzcd}[column sep=large]
 U \arrow[r,"\gamma"] \arrow[d,"g"'] &
 U\boxtimes U \arrow[d,"f_g"]\\
 U \arrow[r,"\delta"'] & U.
\end{tikzcd}
\]
Pullback along $\gamma$ gives
$h^*p_+=h^*d_-$ and $h^*p_-=h^*d_+$.
The separator inequalities imply both
$h^*d_-\le h^*d_+$ and $h^*d_+\le h^*d_-$.
These two domains are complementary, because $h$ is total and
$sh=h$. Equality therefore forces $0_U=1_U$, a contradiction.
\end{proof}

For $p_+=b$ and $p_-=\nu^*b$, every consistent domain extension
$b\le t\le s$ is consequently incomplete: a complete extension would
itself be a relative complementary separator. The representation
condition \eqref{ct:representation} supplies no uniform assignment
of $g$ to a code of $(e_+,e_-)$. Its conclusion is inseparability, as
defined above; effective inseparability requires further uniform data.

\section{The numerical category and effective unions}
\label{num:ambient}

Let $\R_{\mathrm{rec}}$ have objects $\N^k$, for $k\ge1$, and partial
recursive maps as arrows, with equality of partial functions. Composition
is composition of partial maps. Restriction sends a map to the partial
identity on its domain. The restriction product is the Cartesian product
of underlying sets, acting on pairs of partial maps where both components
are defined. Recursive pairing gives isomorphisms $\N^k\cong\N$.
A universal partial recursive function supplies a Turing morphism, making
this the numerical recursion category of \cite{DPH,Heller}.

A set is c.e. precisely when it is the domain of a partial recursive
function. Its partial identity is therefore an arrow. Conversely, the
domain of every restriction idempotent is c.e. For a c.e. subset
$A\subseteq\N^k$, write $1_A$ for this partial identity. The range of
$f$ is $1_{\operatorname{im}f}$, since the image of a partial recursive
map is c.e. It has the universal equality property
\eqref{ct:universal-range}: agreement after $f$ is agreement, including
definedness, at every point of its image.

\begin{proposition}\label{num:uniform-unions}
Suppose $(A_j)_{j\in\N}$ is a uniformly c.e. family of subsets of
$\N^k$, meaning that $\{(j,x):x\in A_j\}$ is c.e. Then
$1_{\bigcup_j A_j}$ is the stable union of the domains $1_{A_j}$ in
$\R_{\mathrm{rec}}$.
\end{proposition}
\begin{proof}
Dovetailing the enumerations gives a c.e. union. If two partial maps agree
after every $1_{A_j}$, they have the same definedness and value at each
point of the union. This proves Heller's universal agreement condition.
For every $m\ge1$, the same argument applies to
$\N^m\times\bigcup_j A_j=\bigcup_j(\N^m\times A_j)$, proving stability.
\end{proof}

The uniformity condition is substantive. Let $B\subseteq\N$ be a set
that is not c.e., and put $B_j=B\cap\{0,\ldots,j\}$. Each $B_j$ is
finite, but the increasing family has no supremum among c.e. sets.
Indeed, a supremum would contain $B$, while every
$\N\setminus\{x\}$ with $x\notin B$ is an upper bound; the supremum
would therefore be $B$.


\section{A numerical deductive construction}
\label{num:section}

An inseparable pair supplies two families of finite computational
codes. Classical deduction makes the families incompatible and permits
extensions to decide assertions for which neither kind of code is
available. Inseparability prevents any consistent effective extension from
deciding every assertion.

\subsection{Computational codes and axioms}

Disjoint c.e. sets $A,B\subseteq\N$ are \emph{recursively inseparable} if no recursive set $C\subseteq\N$
satisfies $A\subseteq C$ and $B\subseteq\N\setminus C$.
Such pairs exist in the numerical category. To see this, fix an effective
numbering $(\varphi_e)_{e\in\N}$ of the unary partial recursive functions
and put
\[
 A=\{e:\varphi_e(e)=0\},\qquad B=\{e:\varphi_e(e)=1\}.
\]
These sets are disjoint and c.e. If a recursive set $C$ separated them,
the total recursive function taking value $1$ on $C$ and $0$ outside $C$
would have an index $d$. If $d\in C$, then $d\in B$, contradicting
separation; if $d\notin C$, then $d\in A$, again contradicting separation.

Fix any disjoint recursively inseparable c.e. pair $A,B$.
Dovetail procedures recognizing membership in $A$ and $B$.
This defines a partial recursive function
\[
 k:\N\rightharpoonup\N,
 \qquad
 k(n)=
 \begin{cases}
 0,&n\in A,\\
 1,&n\in B,
 \end{cases}
\]
which is undefined elsewhere. Disjointness makes the two possible outputs
compatible with determinism.

Fix a machine computing $k$. A \emph{computation code} for
$k(n)=i$ is a finite list of configurations beginning with its initial
configuration on input $n$, following its transition rules, and ending in
a halting configuration with output $i$. Verification of a proposed
code is decidable. Such a code exists exactly when $k(n)=i$.

Let $\mathsf F$ be the set of formulas generated by
\[
 \alpha ::= \bot\mid\top\mid P_n\mid Q_n\mid
 \neg\alpha\mid(\alpha\land\alpha)\mid(\alpha\lor\alpha),
 \qquad n\in\N.
\]
Here $P_n$ and $Q_n$ are distinct propositional atoms. Write
$\alpha\Rightarrow\beta$ for $\neg\alpha\lor\beta$.
Choose an effective coding of finite syntax trees. In particular,
well-formedness and the constructor operations on formula codes are
recursive. 

Define the set of nonlogical axioms $\Gamma_0\subseteq\mathsf F$ by
\begin{equation}\label{num:axioms}
 \Gamma_0=
 \{P_n:n\in A\}\ \cup\ \{Q_n:n\in B\}\ \cup\
 \{\neg(P_n\land Q_n):n\in\N\}.
\end{equation}
The first two families are enumerated by computation codes; the
third is recursive. Thus $\Gamma_0$ is c.e.

For $\Gamma\subseteq\mathsf F$, define $\Gamma\vdash\alpha$ using the
following finite proofs. Each line is either a member of $\Gamma$, a
classical propositional tautology, or the conclusion of modus ponens from
two preceding lines. Whether a formula is a tautology is decidable by
evaluating it on the finitely many truth assignments to its atoms. Put
\[
 \operatorname{Cn}(\Gamma)
   =\{\alpha\in\mathsf F:\Gamma\vdash\alpha\},
 \qquad T_0=\operatorname{Cn}(\Gamma_0).
\]
A \emph{theory} is a set of formulas containing all propositional tautologies
and closed under modus ponens. An \emph{extension} of $T$ is a
theory $T'$ with $T\subseteq T'$. A theory is \emph{consistent} if
$\bot\notin T$, and \emph{complete} if, for every $\alpha\in\mathsf F$,
either $\alpha\in T$ or $\neg\alpha\in T$.

\begin{lemma}\label{num:effective-proofs}
The theory $T_0$ is c.e. It has a presentation by finite proof codes
whose validity is decidable. More generally, the deductive closure of any
c.e. set of axioms is c.e. and has such a code presentation.
\end{lemma}
\begin{proof}
In a proof from $\Gamma_0$, annotate a line $P_n$ or $Q_n$ used as an
axiom with a computation code for its corresponding output. An
exclusion axiom in the third family of \eqref{num:axioms} is recognized
directly. A tautology line is checked by a truth table; a modus ponens line
records the indices of its premises. All checks terminate.

For a general c.e. axiom set, fix an enumerator and annotate each axiom
line with a finite enumeration stage at which it appears. Simulating that
stage is a terminating test. Enumerating all finite annotated proofs and
retaining the valid ones enumerates exactly the deductive closure.
\end{proof}

\begin{lemma}\label{num:consistent}
The theory $T_0$ is consistent.
\end{lemma}
\begin{proof}
Assign $P_n$ the truth value of $n\in A$ and $Q_n$ the truth value of
$n\in B$, and extend this assignment by the Boolean operations. Every
axiom of $\Gamma_0$ is true: the exclusion axioms use $A\cap B=\varnothing$.
Tautologies are true and modus ponens preserves truth. Thus $\bot$ has no
proof from $\Gamma_0$.
\end{proof}

\subsection{Deductive categories and essential incompleteness}

For a theory $T$, define the category $\mathcal B_T$ as follows. Its objects
are formulas in $\mathsf F$, and its hom sets are
\[
 \mathcal B_T(\alpha,\beta)=
 \begin{cases}
 \{*\},&\alpha\Rightarrow\beta\in T,\\
 \varnothing,&\text{otherwise}.
 \end{cases}
\]
The identity exists because $\alpha\Rightarrow\alpha$ is a tautology.
The transitivity tautology and modus ponens give composition. Since every
hom set has at most one member, the identity and associativity laws follow.
Conjunction is a categorical product, disjunction is a coproduct, $\top$
is terminal, and $\bot$ is initial, by the corresponding propositional
tautologies. Negation defines a contravariant functor
\[
 N_T:\mathcal B_T^{\op}\longrightarrow\mathcal B_T
\]
that sends $\alpha$ to $\neg\alpha$ and an arrow
$a^{\op}:\beta\to\alpha$ of $\mathcal B_T^{\op}$ to the arrow
$\neg\beta\to\neg\alpha$ of $\mathcal B_T$.
Contraposition supplies the latter arrow. Double negation gives a natural
isomorphism $1_{\mathcal B_T}\cong N_T\circ N_T^{\op}$. Quotienting formulas by mutual
implication gives the Lindenbaum Boolean algebra of $T$.

For an extension $T\subseteq T'$, define
\[
 J_{T,T'}:\mathcal B_T\longrightarrow\mathcal B_{T'}
\]
to be the identity on formula objects and to send each existing arrow to
the unique arrow with the same endpoints. The inclusion of theorem sets
makes this well-defined. These functors preserve the indicated Boolean
operations and satisfy
$J_{T',T''}J_{T,T'}=J_{T,T''}$ and $J_{T,T}=1$.
Consistency says that $\mathcal B_T(\top,\bot)$ is empty; completeness
says that for each $\alpha$, at least one of
$\mathcal B_T(\top,\alpha)$ and
$\mathcal B_T(\top,\neg\alpha)$ is inhabited.

\begin{theorem}\label{num:essential}
Let $A,B\subseteq\N$ be disjoint recursively inseparable c.e. sets, and
let $T_0$ be the theory defined by \eqref{num:axioms}. Then $T_0$ is a
consistent c.e. theory, and every consistent c.e. extension $T$ of $T_0$
is incomplete. In fact, for some $n$, neither $P_n$ nor $\neg P_n$ belongs
to $T$.
\end{theorem}
\begin{proof}
The preceding lemmas establish the assertions about $T_0$. For an
extension $T$, put
\[
 U_T=\{n:P_n\in T\},\qquad
 V_T=\{n:\neg P_n\in T\}.
\]
Both sets are c.e., since the maps $n\mapsto P_n$ and
$n\mapsto\neg P_n$ on codes are recursive. Consistency implies
$U_T\cap V_T=\varnothing$: from a formula and its negation, the calculus
derives $\bot$.

If $n\in A$, then $P_n\in\Gamma_0\subseteq T$, so $A\subseteq U_T$.
If $n\in B$, then $Q_n$ and $\neg(P_n\land Q_n)$ belong to $T$.
The tautology
\[
 Q_n\Rightarrow
 \bigl(\neg(P_n\land Q_n)\Rightarrow\neg P_n\bigr)
\]
and two applications of modus ponens show that $B\subseteq V_T$.
If $U_T\cup V_T=\N$, their disjoint c.e. enumerations decide membership
in $U_T$: run both enumerations until the input appears. The resulting
recursive set separates $A$ from $B$, contrary to their inseparability.
Hence some $n$ belongs to neither set.
\end{proof}

The consistency argument for $T_0$ uses a truth assignment. The argument
for its extensions uses only consistency, enumerability, and inclusion of
the base axioms. An extension may prove assertions that are false under
that truth assignment.

\subsection{Hoare interpretation and proof domains}

Identify a numerical domain with its c.e. underlying set. Put
$p_n=\{n\}$ and $q_i=\{i\}$ for $i=0,1$. Interpret $P_n$ as the
assertion $k(n)=0$ and $Q_n$ as the assertion $k(n)=1$. With the total- and
partial-correctness profunctors defined above,
\begin{align}
 P_n\text{ is true}
 &\quad\Longleftrightarrow\quad H_k(p_n,q_0)\ne\varnothing,
 \label{num:hoare-positive}\\
 \neg P_n\text{ is true}
 &\quad\Longleftrightarrow\quad K_k(p_n,q_1)\ne\varnothing.
 \label{num:hoare-negative}
\end{align}
Indeed, the first relation requires termination with output $0$. The
second says that every terminating execution from $n$ ends with output
$1$; because $k$ is deterministic and has no other possible outputs, it
holds precisely when $k(n)$ is undefined or equals $1$. These are exactly
the cases in which $P_n$ is false.

The singleton precondition is essential to
\eqref{num:hoare-negative}. For a larger precondition, negating a total
correctness assertion says that at least one input fails its requirement;
a partial correctness assertion constrains every terminating input.
For example, let $f(0)=0$ and $f(1)=1$. Then
$H_f(\{0,1\},\{0\})$ is empty, but
$K_f(\{0,1\},\{1\})$ is empty as well.

Let $s:\N\rightharpoonup\N$ be the partial identity on the set of
well-formed formula codes, and define
\[
 \nu:\N\rightharpoonup\N,\qquad
 \nu(\ulcorner\alpha\urcorner)=\ulcorner\neg\alpha\urcorner
\]
on that set. Then $\bar\nu=s$ and $s\nu=\nu$.
Using Lemma~\ref{num:effective-proofs}, define the partial recursive
conclusion map
\[
 d_0:\N\rightharpoonup\N
\]
to return the last formula of a valid proof code for $T_0$ and to
be undefined on all other inputs. Its range is the partial identity
$t_0$ on the theorem codes of $T_0$, and $t_0\le s$.

Let $a$ and $b$ be the partial identities on $A$ and $B$, respectively,
and define the total recursive map
\[
 \theta:\N\longrightarrow\N,\qquad
 \theta(n)=\ulcorner P_n\urcorner.
\]
There are partial recursive maps $u,v:\N\rightharpoonup\N$ satisfying
\begin{equation}\label{num:certificate-maps}
 \bar u=a,\qquad \bar v=b,\qquad
 d_0u=\theta a,\qquad d_0v=\nu\theta b.
\end{equation}
To construct $u$, simulate $k(n)$ until it returns $0$, and then output
the one-line proof of $P_n$ annotated with the computation just found.
If it returns $1$, continue forever. The construction of $v$ waits instead
for output $1$, and produces the finite proof of $\neg P_n$ used in
Theorem~\ref{num:essential}, with its $Q_n$ line annotated by the trace.
Thus the domains of $u$ and $v$ are exactly $A$ and $B$, and the equations
follow by inspection of their proof conclusions.

For any c.e. extension $T$, its theorem codes similarly form the range
$t$ of a partial recursive proof-conclusion map. The inverse image
$\nu^*t$ consists of the codes of formulas refutable in $T$.
Consistency and completeness become, respectively,
\[
 t\wedge\nu^*t=0,
 \qquad
 t\vee\nu^*t=s.
\]
Consequently \eqref{num:certificate-maps} realizes the maps producing proofs
in the categorical separator argument. Negation is recursive as an
operation on formula codes even though the set of inputs on which
$\neg P_n$ is true need not be c.e.

\subsection{Arithmetic representation of recursive sets}

There is a second arithmetic application, using the representation
hypothesis of Theorem~\ref{ct:representation-inseparable}. Fix an
effective first-order language of arithmetic, an effective classical
deductive calculus, and a recursive coding of finite terms and formulas.
The codes of sentences form a recursive set; let $s$ be its partial
identity. Define $\nu$ on sentence codes by
$\nu(\ulcorner\alpha\urcorner)=\ulcorner\neg\alpha\urcorner$.

Choose a total recursive enumeration
$u\mapsto\varphi_u(v)$ of formulas whose only possible free variable
is $v$, covering every such formula. One construction takes $u$ as a
syntax code when it satisfies the stated condition, and otherwise returns
the fixed formula $0=0$. Write $\underline n$ for the arithmetic numeral
denoting $n$. Numeral substitution defines a total recursive map
\[
 \sigma:\N^2\longrightarrow\N,
 \qquad
 \sigma(u,n)=\ulcorner\varphi_u(\underline n)\urcorner.
\]
The recursion on $n$ constructs the numeral code from the zero and
successor constructors. A finite traversal of the formula tree then
replaces exactly the free occurrences of $v$. Since the replacement term
is closed, it introduces no variable capture. This gives a terminating
algorithm jointly in the formula index $u$ and the numeral $n$.
Every output is a sentence code, so $s\sigma=\sigma$.

A theory $B$ \emph{defines a set $C\subseteq\N$ two-sidedly numeralwise}
if some formula $\varphi(v)$ satisfies
\[
 n\in C\Longrightarrow B\vdash\varphi(\underline n),
 \qquad
 n\notin C\Longrightarrow B\vdash\neg\varphi(\underline n).
\]

\begin{corollary}\label{num:recursive-representation}
Let $B$ be a consistent arithmetical theory in the effective language
above. Suppose that $B$ defines every recursive subset of $\N$
two-sidedly numeralwise. Then every consistent c.e. extension $T$ of
$B$ is incomplete. No enumerability hypothesis on $B$ is required.
\end{corollary}
\begin{proof}
Let $t$ be the partial identity on the theorem codes of $T$, and put
$t^-=\nu^*t$. The domains $t,t^-$ are disjoint by consistency. A
complementary pair of c.e. domains on $\N$ is the pair of partial
identities on a recursive set $C$ and its complement: their enumerations
give a decision procedure by dovetailing.

Choose an index $u$ of a formula defining $C$ two-sidedly numeralwise
in $B$, and define the total constant map $g:\N\to\N$ by $g(x)=u$.
Both families of proofs remain available in $T$. Consistency rules out
the opposite instances, and hence
\[
 [\sigma(g\times1)]^*t=\pi_2^*1_C,
 \qquad
 [\sigma(g\times1)]^*t^-=\pi_2^*1_{\N\setminus C}.
\]
Here $\pi_2:\N^2\to\N$ is the second projection, and the product is
the numerical restriction product. Thus $t,t^-$ satisfy the
representation hypothesis of
Theorem~\ref{ct:representation-inseparable}. If $T$ were complete,
$(t,t^-)$ would itself be a complementary pair relative to $s$ separating
these domains. Their relative inseparability excludes this possibility.
\end{proof}

Di Paola and Heller list several incompleteness conditions, including
conditions involving definability of recursive sets and conditions
involving recursive inseparability \cite[pp.~594--595]{DPH}.

\subsection{Interpretations and effective verification}

The transfer to an arithmetical theory requires an interpretation of the
deductive construction. The following criterion states the data needed
without choosing an arithmetic formalization of the computation relation.

\begin{proposition}\label{num:interpretation}
Let $A,B$ be a disjoint recursively inseparable c.e. pair, and let $S$
be a consistent c.e. classical theory in an effective language.
Suppose $n\mapsto\rho_n$ and $n\mapsto\eta_n$ are computable maps to
codes of sentences of that language, satisfying
\[
 n\in A\Longrightarrow S\vdash\rho_n,
 \qquad
 n\in B\Longrightarrow S\vdash\eta_n,
 \qquad
 S\vdash\neg(\rho_n\land\eta_n)\quad(n\in\N).
\]
Then every consistent c.e. extension of $S$ leaves some $\rho_n$
undecided and is therefore incomplete.
\end{proposition}
\begin{proof}
Send $P_n$ to $\rho_n$, $Q_n$ to $\eta_n$, and extend the assignment by
the Boolean connectives. This is a computable translation $I$ of formulas.
For an extension $S'$, the set
\[
 T_I=\{\alpha\in\mathsf F:S'\vdash I(\alpha)\}
\]
is c.e., contains $\Gamma_0$, and is closed under the propositional
calculus. It is consistent because the translation sends $\bot$ to
$\bot$ and $S'$ is consistent. Theorem~\ref{num:essential} supplies an
$n$ for which neither $P_n$ nor $\neg P_n$ belongs to $T_I$.
\end{proof}

For a proposed interpretation by machine computations in arithmetic,
the hypotheses require provable positive trace instances and provable
incompatibility of the two output assertions. The computability of the
sentence maps includes effective coding and substitution. These are
properties of the chosen interpretation and its base theory.

The effective construction of assertions can be made explicit using finite
computation coding. Given an arithmetic postcondition formula $Q(y)$ and
an arithmetic formula $\operatorname{Trace}_f(n,c,y)$ expressing that $c$
is a halting computation of $f$ from $n$ with output $y$, the formulas
\[
 \exists c\,\exists y\,
   (\operatorname{Trace}_f(n,c,y)\land Q(y)),
 \qquad
 \forall c\,\forall y\,
   (\operatorname{Trace}_f(n,c,y)\Rightarrow Q(y))
\]
define $\operatorname{wp}_f(Q)$ and $\operatorname{wlp}_f(Q)$ in the
standard natural numbers. Bound variables are renamed effectively to avoid
capture. The first formula asks for a successful computation; the second
constrains every possible successful computation. Equation~\eqref{hp:halting-predicates}
shows why this syntactic construction cannot decide membership in the
resulting predicates.

Cook's incompleteness argument for effective program verification uses
the same phenomenon: the partial-correctness assertion
$\{\mathrm{true}\}\,C\,\{\mathrm{false}\}$ holds exactly when $C$
terminates on no initial state \cite[\S6, p.~85]{Cook}. His relative
completeness theorem assumes an expressive assertion language and a
semantically complete assertion proof system \cite[\S6, p.~86]{Cook}.
Theorem~\ref{num:essential} instead treats arbitrary consistent effective
extensions of the displayed deductive system; its proof uses the two
inseparable families of codes.

\section{The intrinsic construction problem}
The transfer theorem applies to a deductive presentation when an inseparable
pair supplies finite proofs of incompatible assertions. The numerical
construction provides the assertion codes, proof codes, and conclusion maps
explicitly. The representation theorem gives a different sufficient
condition: it derives inseparability from the representation equations.
Neither existence of a consequence operator nor its fixed-point
characterization of theories constructs these connecting data.

\begin{problem}\label{prob:intrinsic}
For a recursion category in the sense of Di Paola and Heller, specify a
completion or associated category in which finite computation evidence,
assertions, negation, and deduction are constructed functorially. Prove that
the construction supplies the hypotheses of the code separator
theorem and that its numerical interpretation recovers the stated notion
of consistency and effective extension.
\end{problem}

A solution must explain how the categorical inseparable pair interacts with
deduction. Under the hypotheses of their Theorem~8.15, Di Paola and Heller construct
inseparable domains using a Turing morphism \cite[Definition~8.14 and Theorem~8.15]{DPH}. A Turing morphism is a
universal evaluator: programs for a family of partial maps are supplied by
total parameter maps, and evaluation recovers the family. The role of that
universal computation in an associated deductive category must be specified,
as must the construction of finite evidence and its conclusion maps.

Negation imposes a distinct requirement. The complement of a computational
domain need not itself be a computational domain. Logical negation can
nevertheless be an effective operation on assertion codes, with proof and
refutation searches having computational domains. This is exactly the
separation used by the numerical construction. A completion that adds
complements must also distinguish the original computational domains from
the added predicates: in a Boolean algebra every disjoint pair is separated
by a complemented element, so unrestricted Boolean separation would destroy
the inseparability condition.

The same distinction governs the choice of a semantic target with duality.
An order-reversing involution on assertions describes negation only after its
compatibility with entailment and computational interpretation has been
given. A duality on an arbitrary target category does not supply finite
proofs, their enumeration, or the maps producing proofs. Stable unions assemble
specified computational stages; they do not by themselves construct these
stages or their codes.

An arithmetic interpretation of the numerical deductive system is a further
comparison problem. It requires effective translations of assertions and
proofs that preserve the positive and negative computation codes.
Corollary~\ref{num:recursive-representation} uses an explicitly stated
numeralwise representation hypothesis. Identifying which arithmetic theories
satisfy a particular translation requires a proof in the chosen arithmetic
calculus.


\appendix
\section{Closure operators associated with deductive presentations}
\label{ct:closure-section}

Fix a range restriction category with the restriction products, universal
range property, and finite stable unions of Section~\ref{ct:union-section}.
A specified family of partial rule maps determines a closure operator
when its deduction stages admit stable unions. This describes theories
relative to the supplied rules and axioms.

Fix domains $s\le q$ on an object $U$, and write
$E_q=\{a\in\D(U):a\le q\}$ and $E_s=\{a\in\D(U):a\le s\}$.
The domain $q$ supports all intermediate conclusions; $s$ selects those
conclusions to be regarded as assertions. Fix a domain $\ell\le q$ of
axioms and finitely many rules. A rule with $n_i\ge1$ premises and
parameter object $P_i$ is an arrow
\[
 r_i:P_i\boxtimes U^{\boxtimes n_i}\longrightarrow U,
 \qquad r_i=q r_i(1_{P_i}\boxtimes q^{\boxtimes n_i}).
\]
For a rule without parameters, omit the factor $P_i$. Rules with no
premises are incorporated in $\ell$. Define monotone maps
\[
 R_i:E_q\longrightarrow E_q,\qquad
 R_i(a)=\ran\bigl(r_i(1_{P_i}\boxtimes a^{\boxtimes n_i})\bigr).
\]
A domain $t\le q$ is \emph{rule-closed} when $R_i(t)\le t$ for every $i$.
For $a\le q$, form the increasing deduction sequence
\begin{equation}\label{ct:stages}
 a_0=\ell\vee a,\qquad
 a_{k+1}=a_k\vee\bigvee_i R_i(a_k).
\end{equation}

\begin{theorem}[Closure under finite deduction]\label{ct:closure}
Suppose every sequence \eqref{ct:stages} has a stable union. The functor
$c_q:E_q\to E_q$ defined by $c_q(a)=\bigvee_k a_k$ and its monotone
action on inequalities is a closure operator. Its value is the least
rule-closed domain containing $\ell\vee a$.
\end{theorem}
\begin{proof}
By Lemma~\ref{ct:union-calculus},
\[
 R_i(c_q(a))=\bigvee_k R_i(a_k)
 \le\bigvee_k a_{k+1}=c_q(a).
\]
Hence $c_q(a)$ is rule-closed. If $t$ is rule-closed and contains
$\ell\vee a$, induction gives $a_k\le t$ for every $k$, so $c_q(a)\le t$.
Monotonicity follows by comparing the stages. Extensivity follows from
$a\le a_0$, and leastness gives $c_q^2=c_q$.
\end{proof}

The selected closure is the functor $c_s:E_s\to E_s$ with
$c_s(a)=s c_q(a)$. It is extensive and monotone; moreover,
\[
 c_s(c_s(a))=s c_q(s c_q(a))\le s c_q(c_q(a))=c_s(a),
\]
so it is idempotent. Let $\Th_c$ be the full subposet of $E_s$ on the
fixed points of $c_s$. Its objects are called \emph{theories}; an
\emph{extension} of $b$ is a theory $t$ with $b\le t$.
Define functors
\[
 L:E_s\longrightarrow\Th_c,
 \qquad j:\Th_c\longrightarrow E_s
\]
by $L(a)=c_s(a)$ and by inclusion, respectively, with their induced
actions on inequalities. Then $L\dashv j$, since for a closed $t$,
\[
 c_s(a)\le t\quad\Longleftrightarrow\quad a\le t.
\]

\subsection{Presentations and the natural inclusion of theories}

A presentation $\mathcal P$ consists of the data
$(U,q,s,\ell,(r_i))$ above, satisfying the deduction-chain hypothesis.
Define the category $\Cons$ to have these presentations as objects.
A morphism $f:\mathcal P\to\mathcal P'$ is an arrow $f:U\to U'$ with
\begin{equation}\label{ct:continuous}
 \bar f=s,\qquad s'f=f,\qquad
 c_s f^*(a)\le f^*c_{s'}(a)\quad(a\le s').
\end{equation}
The first two equations say that $f$ is total from the selected domain
$(U,s)$ to $(U',s')$ in the category obtained by splitting restriction
idempotents: arrows there satisfy $s'f=f=fs$, and the identity of
$(U,s)$ is $s$. Thus the identity in $\Cons$ is $s$, and composition
is composition in $\R$. For composable $f,g$, the support equations
give $\overline{gf}=s$, while monotonicity gives
$c_s f^*g^*\le f^*c_{s'}g^*\le f^*g^*c_{s''}$.

Define contravariant functors
\[
 \mathcal E,\Th:\Cons^{\op}\longrightarrow\Pos.
\]
On a presentation they assign $\mathcal E(\mathcal P)=E_s$ and
$\Th(\mathcal P)=\Th_c$. On $f^{\op}:\mathcal P'\to\mathcal P$,
both act by $f^*$, restricted to the indicated posets. If $t$ is closed,
\eqref{ct:continuous} and extensivity imply $c_s(f^*t)=f^*t$;
therefore the latter restriction is defined. The functor laws follow
from those of inverse image.

The inclusions define a natural transformation $j:\Th\Rightarrow\mathcal E$.
Its naturality square is
\[
\begin{tikzcd}[column sep=large]
 \Th(\mathcal P') \arrow[r,"j_{\mathcal P'}"]
   \arrow[d,"\Th(f^{\op})"'] &
 \mathcal E(\mathcal P') \arrow[d,"\mathcal E(f^{\op})"]\\
 \Th(\mathcal P) \arrow[r,"j_{\mathcal P}"'] & \mathcal E(\mathcal P).
\end{tikzcd}
\]
Both paths send $t$ to $f^*t$. Each component is an order embedding,
so $j$ is a monomorphism in the functor category
$[\Cons^{\op},\Pos]$. The adjunctions $L\dashv j$ are pointwise;
naturality of $L$ would require equality in \eqref{ct:continuous}.

The target $\Pos$ is necessary for these general closure data.
For example, on $\mathcal P(\{a,b,z\})$, let the closure operator add $z$
exactly when both $a$ and $b$ are present. The closed sets $\{a\}$ and
$\{b\}$ have closed join $\{a,b,z\}$, whereas their ambient join is
$\{a,b\}$. Thus inclusion of fixed points need not preserve joins.


\subsection{Numerical deduction stages}

For a finite family of partial recursive rule maps and a c.e. axiom
domain, the stages in \eqref{ct:stages} are uniformly c.e. From indices
for the preceding stage one can compute indices for its finite products,
their images under the rule maps, and the finite union of those images.
Iteration gives a stage index uniformly in the stage number.
Proposition~\ref{num:uniform-unions} thus supplies the deduction-chain
hypothesis of Theorem~\ref{ct:closure}.


The closure construction also has a direct realization. Take $U=\N$,
$q=s$ equal to the formula domain, $\ell$ equal to the domain of the
nonlogical axioms and propositional tautologies, and take modus ponens
as the partial recursive binary rule. The stable union of its uniformly
c.e. stages consists exactly of conclusions of finite proofs: induction
on stages gives proofs, and induction on proof length bounds the stage
at which a conclusion appears. Thus the abstract closure is
$A\mapsto\operatorname{Cn}(\Gamma_0\cup A)$ on c.e. formula sets;
its fixed points are the c.e. deductively closed extensions of $T_0$.

For any set $\Gamma$ of formulas, the operator $\operatorname{Cn}$
is extensive and monotone. Replacing uses of derived formulas by their
finite proofs proves idempotence. Its fixed points are exactly the sets
containing all tautologies and closed under modus ponens, that is, the
theories of Section~\ref{num:section}. For a c.e. axiom set, its deductive
closure is a computational domain by Proposition~\ref{num:uniform-unions}.

\section{Ordinary interpretations and finite hom-sets}
\label{al:finite}

An interpretation of computations by an ordinary functor can lose all
distinctions between parallel programs. The following obstruction uses only
composition, identities, and a family of partial retractions.

\begin{proposition}\label{al:collapse}
Let $\mathcal C$ be a category with zero morphisms and an object $U$.
Suppose that there are endomorphisms $u_n,v_n:U\to U$, for $n\in\N$, such that
\[
 v_nu_n=1_U,\qquad v_nu_m=0_{U,U}\quad(n\ne m).
\]
If $F:\mathcal C\to\mathcal B$ is a functor and
$\mathcal B(FU,FU)$ is finite, then $F(1_U)=F(0_{U,U})$.
For objects $X,Y$ isomorphic to $U$, all arrows $X\to Y$ have the same
image under $F$.
\end{proposition}
\begin{proof}
There are distinct $i,j$ with $F(u_i)=F(u_j)$. Composing with $F(v_i)$ gives
\[
 F(1_U)=F(v_i)F(u_i)=F(v_i)F(u_j)=F(0_{U,U}).
\]
If $a:U\to X$ is an isomorphism, then
\[
 F(1_X)=F(a)F(1_U)F(a^{-1})
       =F(a0_{U,U}a^{-1})=F(0_{X,X}).
\]
Consequently, every $f:X\to Y$ satisfies
$F(f)=F(f)F(0_{X,X})=F(0_{X,Y})$.
\end{proof}

A category is \emph{locally finite} if each of its hom-sets is finite.
In the category $\R_{\mathrm{rec}}$ of partial recursive maps between the
objects $\N^k$, $k\ge1$, choose a recursive pairing bijection
$c:\N^2\to\N$ with recursive projections. Define
\[
 u_n(x)=c(n,x),\qquad
 v_n(c(m,x))=
 \begin{cases}x,&m=n,\\ \text{undefined},&m\ne n.\end{cases}
\]
These are partial recursive maps with the equations of
Proposition~\ref{al:collapse}. Recursive coding also gives isomorphisms
$\N^k\cong\N$. Hence every functor from $\R_{\mathrm{rec}}$ to a locally
finite category identifies all parallel arrows. No preservation of zero
morphisms is required. In particular, a faithful ordinary interpretation
of $\R_{\mathrm{rec}}$ must give $\N$ an image with infinitely many
endomorphisms.

\section{Adjoining partial witnesses to a profunctor}
\label{al:extension}

An interpretation already defined for total maps can be extended by
adjoining partial witnesses. Here the original interpretation is recovered
as the part with total support.

Let $\R$ be a small restriction category, let $\mathcal C=\Tot(\R)$ be its
wide subcategory of total maps, and let $J:\mathcal C\to\R$ be the inclusion.
Let $\mathcal B$ be a small category, and suppose that a profunctor
\[
 Q:\mathcal C^{\op}\times\mathcal B\longrightarrow\Set
\]
is given. Define a functor
\[
 P:\R^{\op}\times\mathcal B\longrightarrow\Set
\]
on objects by the coend
\begin{equation}\label{al:coend}
 P(X,A)=\int^{Y\in\mathcal C}\R(X,Y)\times Q(Y,A).
\end{equation}
Explicitly, $P(X,A)$ is the quotient of the disjoint union of these products
by the equivalence relation generated by
\begin{equation}\label{al:relation}
 [gh,q']=[h,Q(g^{\op},1_A)(q')],
\end{equation}
where $h:X\to Y$ is in $\R$, $g:Y\to Y'$ is total, and
$q'\in Q(Y',A)$. The symbol $[h,q]$ denotes the resulting class.
For $f:X'\to X$ in $\R$ and $\alpha:A\to A'$ in $\mathcal B$, put
\begin{equation}\label{al:action}
 P(f^{\op},\alpha)[h,q]
   =[hf,Q(1_Y^{\op},\alpha)(q)].
\end{equation}
The functor laws of $Q$ show that this assignment preserves
\eqref{al:relation}; associativity of composition and the same functor laws
then prove the identity and composition laws for $P$.

\begin{proposition}\label{al:support}
For each $X\in\R$ and $A\in\mathcal B$, there is a well-defined map
\[
 \delta_{X,A}:P(X,A)\longrightarrow\Dom(X),\qquad
 \delta_{X,A}[h,q]=\bar h.
\]
It satisfies
\[
 \delta_{X',A'}\bigl(P(f^{\op},\alpha)(p)\bigr)
   =f^*\bigl(\delta_{X,A}(p)\bigr).
\]
The maps $q\mapsto[1_X,q]$ identify $Q(X,A)$ with the fibre
$\delta_{X,A}^{-1}(1_X)$, naturally in $X\in\mathcal C$ and
$A\in\mathcal B$.
\end{proposition}
\begin{proof}
For total $g$, the restriction identity gives
$\overline{gh}=\bar h$. Thus both sides of each generating relation
\eqref{al:relation} have the same restriction, proving that $\delta$ is
well-defined. The restriction identities also give
\[
 \overline{hf}=\overline{\bar h f}=f^*(\bar h),
\]
which proves the displayed compatibility with $P$.

The class $[1_X,q]$ has support $1_X$. Conversely, if $[h,q]$ has support
$1_X$, then $h$ is total and \eqref{al:relation} gives
\[
 [h,q]=[1_X,Q(h^{\op},1_A)(q)].
\]
On representatives with total support the assignment
$[h,q]\mapsto Q(h^{\op},1_A)(q)$ respects the generating relation, by
functoriality of $Q$. Every representative of such a class has total
support, so this assignment defines the inverse to $q\mapsto[1_X,q]$.
For total $f$, substituting \eqref{al:action} and then
\eqref{al:relation} verifies naturality.
\end{proof}

Thus there is a componentwise injective natural transformation
\[
 Q\longrightarrow P\circ(J^{\op}\times1_{\mathcal B}).
\]
The coend specifies how partial witnesses transform once $Q$ has been
supplied. Its input still includes the objects, morphism action, and
logical interpretation of $Q$. Moreover, an order-reversing operation on
$\mathcal B$ would change the variance of that argument; it does not by
itself induce an operation on the sets $P(X,A)$ of the same variance.

\section{Indexed predicates and enrichment}
\label{al:enrichment}

The domain functor $\D:\R^{\op}\to\Pos$ is an indexed poset. Its objects
$\D(X)$ are posets of predicates, and its morphism maps reindex predicates
along computations. A natural inclusion between two such functors is a
morphism in an ordinary functor category. An enrichment over a category
of computational objects requires further data.

In Kelly's terminology, enrichment requires a monoidal category
\[
 \mathscr V=(\mathscr V_0,\otimes,a,I,\ell,r).
\]
Here $\otimes:\mathscr V_0\times\mathscr V_0\to\mathscr V_0$ is a functor,
$I$ is its unit object, and $a,\ell,r$ are coherent associativity and unit
isomorphisms. A $\mathscr V$-category $\mathcal A$ has objects and, for each
pair $X,Y$, an object $\mathcal A(X,Y)$ of $\mathscr V_0$, together with
morphisms in $\mathscr V_0$
\[
 \mathcal A(Y,Z)\otimes\mathcal A(X,Y)\longrightarrow\mathcal A(X,Z),
 \qquad I\longrightarrow\mathcal A(X,X),
\]
satisfying associativity and unit laws \cite[\S\S1.1--1.2]{Kelly}.
For example, with $\mathscr V_0=\mathbf{Ab}$ and tensor product
$\otimes_{\mathbb Z}$, every hom-object is an abelian group and composition
is bilinear.

The Hoare constructions use ordinary profunctors with values in $\Set$.
Their dependence on recursion-theoretic structure occurs in the domain
posets, inverse images, ranges, and effective numerical presentations.
An additional enrichment claim would require a specified monoidal category
and hom-objects with the displayed composition and unit morphisms.

Logical negation also has a separate role from reindexing. If
$\nu:U\to U$ is a partial map coding negation of assertions, its action
on domains is the monotone map
\[
 \nu^*:\Dom(U)\longrightarrow\Dom(U),\qquad
 t\longmapsto\overline{t\nu}.
\]
This inverse image selects assertions whose negations belong to $t$.
By contrast, Boolean negation reverses the order of logical implication.
In the numerical model, set complementation cannot define an operation on
all restriction domains: the complement of a nonrecursive c.e. subset of
$\N$ is not c.e. A category of logical assertions can therefore have
negation even when its computational interpretation requires predicates
outside the poset of restriction domains.

\begin{thebibliography}{99}
\bibitem{DPH} R.~A.~Di Paola and A.~Heller,
Dominical categories: recursion theory without elements,
\emph{Journal of Symbolic Logic} \textbf{52}(3) (1987), 594--635.
\href{https://doi.org/10.2307/2274352}{doi:10.2307/2274352}.
\bibitem{Heller} A.~Heller,
An existence theorem for recursion categories,
\emph{Journal of Symbolic Logic} \textbf{55}(3) (1990), 1252--1268.
\href{https://doi.org/10.2307/2274486}{doi:10.2307/2274486}.
\bibitem{CockettLack} J.~R.~B.~Cockett and S.~Lack,
Restriction categories I: categories of partial maps,
\emph{Theoretical Computer Science} \textbf{270}(1--2) (2002), 223--259.
\bibitem{CockettGarner} J.~R.~B.~Cockett and R.~Garner,
Restriction categories as enriched categories,
\emph{Theoretical Computer Science} \textbf{523} (2014), 37--55.
\href{https://arxiv.org/abs/1211.6170}{arXiv:1211.6170}.
\bibitem{Kelly} G.~M.~Kelly,
\emph{Basic Concepts of Enriched Category Theory},
London Mathematical Society Lecture Note Series~64,
Cambridge University Press, 1982;
reprinted in \emph{Reprints in Theory and Applications of Categories}
\textbf{10} (2005), 1--136.
\url{https://www.tac.mta.ca/tac/reprints/articles/10/tr10.pdf}.
\bibitem{Dijkstra} E.~W.~Dijkstra,
Guarded commands, non-determinacy and a calculus for the derivation of programs,
EWD418, author manuscript.
\url{https://www.cs.utexas.edu/~EWD/transcriptions/EWD04xx/EWD418.html}.
\bibitem{DijkstraScholten} E.~W.~Dijkstra and C.~S.~Scholten,
\emph{Predicate Calculus and Program Semantics},
Springer, New York, 1990.
\href{https://doi.org/10.1007/978-1-4612-3228-5}{doi:10.1007/978-1-4612-3228-5}.
\bibitem{Cook} S.~A.~Cook,
Soundness and completeness of an axiom system for program verification,
\emph{SIAM Journal on Computing} \textbf{7}(1) (1978), 70--90.
\url{https://www.cs.toronto.edu/~sacook/homepage/soundness.pdf}.
\bibitem{ArthanEtAl} R.~Arthan, U.~Martin, E.~A.~Mathiesen, and P.~Oliva,
A general framework for sound and complete Floyd--Hoare logics,
\emph{ACM Transactions on Computational Logic} \textbf{11}(1) (2009),
Article~7, 1--31.
\href{https://doi.org/10.1145/1614431.1614438}{doi:10.1145/1614431.1614438};
\href{https://arxiv.org/abs/0807.1016}{arXiv:0807.1016}.
\end{thebibliography}
\end{document}
